\section{Angular inversion of the leading massive QED charge}
\label{sec:massive-inverse}

Campiglia and Laddha's leading massive QED Ward identity uses the
hyperboloid extension of the sphere gauge parameter \cite{massive-qed}.
We start from its geometric kernel and analyze what angular charge data
can recover. This section assumes a prescribed signed hard-charge
density and a single known speed. It does not derive a higher
logarithmic Ward identity.

Use the round unit sphere with area $4\pi$ and Laplacian
$\Delta_{S^2}=\operatorname{div}\nabla$. For a velocity $v=\beta u$ with
$u\in S^2$, $0<\beta<1$, the tangent soft factor in signature $(-+++)$
is $\partial_A\log(1-\beta u\cdot n)$: it is the ratio
$p\cdot (0,\partial_A n)/[p\cdot (1,n)]$, with $p=E(1,\beta u)$.
Differentiating in $t=u\cdot n$ gives
\[
 \Delta_{S^2}\log(1-\beta t)
 =\frac{1-\beta^2}{(1-\beta t)^2}-1.
\]
Thus in units where the hard charge multiplies the normalized angular
kernel, the sphere divergence divided by $4\pi$ is
$K_\beta(t)-1/(4\pi)$, with
\[
 K_\beta(t)=\frac{1-\beta^2}{4\pi(1-\beta t)^2}.
\]
This is also the hyperboloid kernel
$[4\pi (Y^0-\mathbf Y\cdot n)^2]^{-1}$ for
$Y=(1,\beta u)/\sqrt{1-\beta^2}$. Direct integration gives
$\int_{S^2} K_\beta(u\cdot n)\,dn=1$. Summing the signed contributions
cancels the constant term when the total signed charge is zero. Equivalently,
the result is restricted to angular modes of positive degree.

Define $(T_\beta\rho)(n)=\int_{S^2}K_\beta(u\cdot n)\rho(u)\,du$.
Let $Y_{\ell j}$ be a real orthonormal spherical harmonic basis and use
the norm $\|\rho\|_{H^s}^2=\sum_{\ell,j}
 (1+\ell(\ell+1))^s|\rho_{\ell j}|^2$.

\begin{theorem}[Sharp logarithmic recovery at a known speed]
\label{thm:massive-inverse}
Put $r=\beta/(1+\sqrt{1-\beta^2})\in(0,1)$.
For constants $0<c_\beta\le C_\beta<\infty$, the eigenvalue $\lambda_\ell$
of $T_\beta$ on degree $\ell$ satisfies
\[
 c_\beta\sqrt{\ell+1}\,r^\ell
 \ \le\ \lambda_\ell\ \le\
 C_\beta\sqrt{\ell+1}\,r^\ell\qquad(\ell\ge0).
\]
More precisely, $\lambda_\ell/(\sqrt{\ell}\,r^\ell)$ tends to
$\sqrt{\pi(1-r^2)}/2$ as $\ell\to\infty$.
In particular, the exact data determine $\rho$. For $s>0$, let
$\mathcal D=\{\rho:\int\rho=0,\ \|\rho\|_{H^s}\le M\}$ and observe
$y=T_\beta\rho+e$ with $\|e\|_2\le\eta$. As $\eta/M\to0$, the infimum
over all reconstruction rules of the worst-case $L^2$ error on
$\mathcal D$ is bounded above and below by positive multiples of
$M[\log(M/\eta)]^{-s}$ depending only on $\beta,s$.
\end{theorem}

\begin{proof}
By Funk--Hecke, $\lambda_\ell=(1-\beta^2)\int_{-1}^1
 P_\ell(t)(1-\beta t)^{-2}\,dt/2$ for $P_\ell(1)=1$.
Rodrigues' formula and $\ell$ integrations by parts give
\[
 \lambda_\ell=(\ell+1)\int_{-1}^1 h(t)^\ell\,w(t)\,dt,\quad
 h(t)=\frac{\beta(1-t^2)}{2(1-\beta t)},\quad
 w(t)=\frac{1-\beta^2}{2(1-\beta t)^2}.
\]
The boundary terms vanish since $(t^2-1)^\ell$ and its derivatives of
order less than $\ell$ vanish at both endpoints. The differentiated
kernel is $(\ell+1)!\,\beta^\ell(1-\beta t)^{-\ell-2}$.
The weight $w$ has integral one and is bounded above and below by
positive constants. Since $\beta=2r/(1+r^2)$,
\[
 1-\frac{h(t)}r=\frac{(t-r)^2}{1+r^2-2rt}.
\]
For $\ell\ge1$, $(h/r)^\ell\le
\exp[-\ell(t-r)^2/(1+r)^2]$. The Gaussian integral on the real line gives
the upper bound. For the lower bound, restrict to
$|t-r|\le(1-r)/(2\sqrt{\ell+1})$ inside $[-1,1]$. There
$(h/r)^\ell\ge (1-1/[4(\ell+1)])^\ell\ge 3/4$ by Bernoulli's
inequality. The interval length gives the lower bound. Finally
$\lambda_0=1$.

To obtain the constant, set $t=r+y/\sqrt{\ell}$ in the positive
integral and extend its rescaled integrand by zero outside the
rescaled interval. The expression $(h(t)/r)^\ell$ tends pointwise to
$\exp(-y^2/(1-r^2))$, and $w(r)=1/2$. The same Gaussian upper
bound and boundedness of $w$ justify dominated convergence.

Invert the modes $1\le\ell\le L$ and set the others to zero.
The squared-error contributions are on disjoint angular modes. In
particular the total error is at most
\[
 M(1+(L+1)(L+2))^{-s/2}+c_\beta^{-1}\eta r^{-L}.
\]
Take $L=\lfloor \log(M/\eta)/(2|\log r|)\rfloor$ for $\eta/M$
sufficiently small. The first term gives the logarithmic bound and
the second is at most $c_\beta^{-1}\sqrt{M\eta}$.
For the lower bound set
$\ell=\lceil 2\log(M/\eta)/|\log r|\rceil$ and take the two profiles
$\rho_\pm=\pm M(1+\ell(\ell+1))^{-s/2}Y_{\ell j}$. Both satisfy the
prior. Their data norms are at most
$C_\beta M\sqrt{\ell+1}\,(\eta/M)^2\le\eta$ for small $\eta/M$.
Zero is then a common possible observation. The triangle inequality
forces error at least $\|\rho_+\|_2$ for one of them, giving the lower bound.
\end{proof}

If speed mixtures are allowed as arbitrary densities without
additional observations, injectivity already fails for two distinct
speeds: a harmonic of positive degree carried at speed $\beta_1$ and
the same harmonic multiplied by $\lambda_\ell(\beta_1)/\lambda_\ell(\beta_2)$
at speed $\beta_2$ give identical data. Both can be scaled to meet any
positive Sobolev bound. This is a limitation of an angular observation
on signed densities. It neither asserts a new scattering identity nor
rules out recovery from extra energy-resolved data or a more restrictive
model of the hard particles.
